% ============================================
% Глава 1. Обзор литературы
% ============================================

\chapter{ОБЗОР ПРЕДМЕТНОЙ ОБЛАСТИ}

\section{Системная динамика: методология и приложения}

\subsection{Исторические предпосылки и фундаментальные основы}

Системная динамика представляет собой методологию моделирования и анализа поведения сложных систем, характеризующихся наличием нелинейных взаимосвязей, обратных связей и запаздываний. Методология была основана Джейом Форрестером в 1950--1960-х годах в Массачусетском технологическом институте. Первоначально разработанная для анализа промышленных систем, методология быстро нашла применение вUrban Dynamics \cite{forrester1969urban}, изучении глобальных проблем \cite{meadows1972limits} и бизнес-стратегиях \cite{sterman2000business}.

Фундаментальные принципы системной динамики были изложены в работе \cite{forrester1961industrial}, где Форрестер ввёл ключевые понятия: \emph{уровни} (stocks, накопители), \emph{потоки} (flows) и \emph{обратные связи} (feedback loops). Уровни представляют собой состояния системы, характеризующиеся наличием памяти --- они аккумулируют значения на протяжении времени. Потоки определяют скорости изменения уровней и могут зависеть от текущих значений уровней, внешних параметров и времени. Обратные связи формируют структуру взаимодействия элементов системы и определяют её поведение во времени.

Поведение системно-динамической модели определяется решением системы обыкновенных дифференциальных уравнений первого порядка, которые в общем виде записываются следующим образом:

\begin{equation}
    \frac{dS_i}{dt} = F_i(S_1, S_2, \ldots, S_n; P_1, P_2, \ldots, P_m; t), \quad i = 1, \ldots, n
    \label{eq:sde_general}
\end{equation}

где $S_i$ --- уровни (state variables) модели, $F_i$ --- потоки, $P_j$ --- внешние параметры, $t$ --- время, $n$ --- число уровней, $m$ --- число параметров.

Форрестер выделил два фундаментальных типа петель обратной связи: \emph{усиливающие} (reinforcing, положительная обратная связь) и \emph{балансирующие} (balancing, отрицательная обратная связь). Усиливающие петли порождают экспоненциальный рост или коллапс --- поведение, при котором небольшие возмущения усиливаются системой. Балансирующие петли, напротив, стремятся привести систему к равновесному состоянию, компенсируя отклонения от целевого значения. Взаимодействие этих двух типов петель определяет основные архетипы поведения системы.

\subsection{Концептуальное моделирование: диаграммы причинно-следственных связей}

Перед формализацией математической модели исследователь выполняет качественный этап концептуального моделирования, результатом которого является \emph{диаграмма причинно-следственных связей} (Causal Loop Diagram, CLD). CLD представляет собой ориентированный граф, вершины которого соответствуют переменным системы, а рёбра --- причинно-следственным влияниям между переменными.

Каждое ребро графа помечается знаком: плюс (+) означает, что увеличение причины ведёт к увеличению следствия (при прочих равных), минус (-) означает, что увеличение причины ведёт к уменьшению следствия. Петля обратной связи классифицируется как усиливающая, если число отрицательных связей в ней чётно, и как балансирующая, если число отрицательных связей нечётно.

CLD служит важнейшим инструментом коммуникации между участниками моделирования, поскольку позволяет представить структуру системы в наглядной форме, доступной для понимания специалистами без математической подготовки. Кроме того, CLD фиксирует гипотезы исследователя о структуре системы и может использоваться для структурного анализа до перехода к количественному моделированию.

\subsection{Диаграммы потоков и уровней}

На следующем этапе моделирования CLD трансформируется в \emph{диаграмму потоков и уровней} (Stock and Flow Diagram, SFD). SFD является более формальным представлением модели, непосредственно отражающим структуру системы дифференциальных уравнений. В SFD вводятся специальные графические элементы:

\begin{itemize}
    \item \textbf{Уровни} (stocks) --- изображаются в виде прямоугольников и представляют накопители, интегрирующие потоки;
    \item \textbf{Потоки} (flows) --- изображаются в виде клапанов (стрелок с регулятором) и определяют скорость изменения уровней;
    \item \textbf{Вспомогательные переменные} (auxiliaries) --- изображаются в виде окружностей и представляют промежуточные вычисления;
    \item \textbf{Константы и экзогенные переменные} --- также изображаются в виде окружностей или ромбов;
    \item \textbf{Источники и стоки} (sources/sinks) --- изображаются в виде облаков и представляют границы моделируемой системы.
\end{itemize}

Каждый уровень в SFD соответствует переменной состояния в системе уравнений, а каждый поток --- правой части дифференциального уравнения. Топология SFD однозначно определяет структуру системы уравнений, хотя конкретные функциональные зависимости должны быть заданы отдельно.

\subsection{Применения системной динамики}

Системная динамика нашла применение в широком спектре предметных областей. Стерман \cite{sterman2000business} систематизировал применение методологии в бизнес-стратегиях, операционном управлении и управлении цепями поставок. В области экологии работы \cite{meadows1972limits} продемонстрировали способность системно-динамических моделей описывать глобальные процессы и выявлять режимы устойчивого развития.

В сфере здравоохранения системная динамика используется для моделирования распространения инфекционных заболеваний, планирования ресурсов медицинских учреждений и оценки вмешательств в системы общественного здравоохранения \cite{rahimi2020healthcare}. В энергетике методология применяется для анализа перехода к возобновляемым источникам энергии и моделирования энергетических рынков \cite{keith2017energy}.

Несмотря на широкий спектр приложений, методология сталкивается с ограничениями, связанными с трудоёмкостью процесса моделирования, необходимостью в специализированной экспертизе и сложностью валидации моделей. Именно эти ограничения определяют потенциальную ценность автоматизации с использованием современных методов искусственного интеллекта.


\section{Программные платформы системной динамики}

\subsection{Коммерческие системы моделирования}

Современный ландшафт инструментов системной динамики включает как коммерческие, так и открытые программные продукты. Наиболее широкое распространение получили следующие системы.

\textbf{Vensim} (Ventana Systems) является одним из наиболее авторитетных инструментов в сообществе системной динамики. Система поддерживает все этапы моделирования --- от построения CLD до симуляции и анализа чувствительности. Vensim реализует собственный язык моделирования и формат хранения моделей (файлы с расширением \texttt{.mdl}), который фактически является стандартом де-факто в области. Расширенные версии (Vensim DSS и Professional) включают инструменты калибровки моделей по данным, оптимизации и анализа чувствительности.

\textbf{Stella} (isee systems) --- один из старейших инструментов системной динамики, изначально разработанный для платформы Macintosh. Stella предлагает визуальный интерфейс моделирования с богатыми возможностями построения диаграмм и интерактивного анализа. Система позиционируется как образовательный инструмент и широко используется в академических учреждениях.

\textbf{AnyLogic} (The AnyLogic Company) представляет собой мультиметодную платформу моделирования, поддерживающую системную динамику наряду с агентным и дискретно-событийным моделированием. AnyLogic обеспечивает интеграцию различных парадигм моделирования в рамках единого проекта, что позволяет создавать гибридные модели сложных систем.

\textbf{Powersim Studio} --- коммерческая платформа, ориентированная на построение моделей для бизнес-анализа и принятия решений. Система включает инструменты разработки интерактивных панелей управления (dashboards) и интеграции с внешними источниками данных.

\subsection{Открытые инструменты и библиотеки}

Важным направлением развития инструментария системной динамики является создание открытых программных решений, обеспечивающих доступность методологии и возможность интеграции с другими системами.

\textbf{PySD} \cite{pysd2021} --- библиотека на языке Python, разработанная Хинкельманом и др., обеспечивающая парсинг, трансляцию и симуляцию моделей системной динамики, созданных в Vensim и Stella. PySD транслирует модели из формата \texttt{.mdl} в код Python, используя библиотеку \texttt{pandas} для представления временных рядов и \texttt{scipy} для численного интегрирования. Ключевые возможности PySD включают:

\begin{itemize}
    \item загрузку моделей Vensim (формат \texttt{.mdl}) и XMILE (формат \texttt{.stmx});
    \item симуляцию моделей с использованием различных методов интегрирования (Эйлера, Рунге-Кутта);
    \item выполнение параметрического анализа и анализа чувствительности;
    \item экспорт моделей в формате Python для интеграции с экосистемой научных вычислений.
\end{itemize}

PySD является ключевым компонентом программной инфраструктуры проекта \texttt{sdCloud}, поскольку обеспечивает возможность симуляции сгенерированных LLM моделей без необходимости использования коммерческого ПО.

\textbf{BPTK-Py} --- фреймворк на Python для системной динамики и агентного моделирования, разработанный компанией transentis. Фреймворк поддерживает моделирование в духе System Dynamics с использованием современных возможностей Python.

\textbf{Mesa} --- библиотека для агентного моделирования на Python, которая может использоваться в сочетании с компонентами системной динамики для создания гибридных моделей.

\subsection{Форматы представления моделей}

Важным аспектом интеграции различных инструментов системной динамики является унификация форматов представления моделей. Формат \textbf{XMILE} (XML Model Interchange Language) \cite{xmile2015} был разработан как стандарт обмена моделями системной динамики на базе XML. XMILE позволяет переносить модели между различными программными продуктами и обеспечивает машинную читаемость структуры модели.

Формат \textbf{MDL} (Vensim Model) использует текстовое представление, где уравнения модели записываются в декларативном стиле. Каждое уравнение имеет вид:

\begin{lstlisting}[language={}, caption={Пример записи уравнения в формате MDL}]
Variable = Expression
~ Units | Commentary
\end{lstlisting}

где \texttt{Expression} может содержать ссылки на другие переменные, функции (INTEG, SMOOTH, DELAY и др.) и константы.

Наличие стандартизированных форматов представления моделей является важным фактором для автоматизированной генерации, поскольку позволяет формализовать целевой формат выходных данных LLM.


\section{Большие языковые модели}

\subsection{Эволюция языковых моделей}

Большие языковые модели (LLM) представляют собой класс нейросетевых архитектур, основанных на механизме внимания (attention mechanism) \cite{vaswani2017attention}, обученных на массивных корпусах текстовых данных для моделирования последовательностей естественного языка. Ключевым событием в развитии LLM стала публикация архитектуры GPT (Generative Pre-trained Transformer) \cite{radford2018gpt} и её последующие итерации: GPT-2 \cite{radford2019gpt2}, GPT-3 \cite{brown2020gpt3} и GPT-4 \cite{openai2023gpt4}.

Архитектура трансформера основана на механизме самовнимания (self-attention), который позволяет модели устанавливать зависимости между элементами входной последовательности независимо от их расстояния. Формально, для входной последовательности $\mathbf{X} = (\mathbf{x}_1, \ldots, \mathbf{x}_n)$ вычисляются запросы (queries), ключи (keys) и значения (values):

\begin{equation}
    \mathbf{Q} = \mathbf{X}\mathbf{W}_Q, \quad \mathbf{K} = \mathbf{X}\mathbf{W}_K, \quad \mathbf{V} = \mathbf{X}\mathbf{W}_V
\end{equation}

где $\mathbf{W}_Q$, $\mathbf{W}_K$, $\mathbf{W}_V$ --- обучаемые матрицы проекции. Веса внимания вычисляются как:

\begin{equation}
    \text{Attention}(\mathbf{Q}, \mathbf{K}, \mathbf{V}) = \text{softmax}\left(\frac{\mathbf{Q}\mathbf{K}^T}{\sqrt{d_k}}\right)\mathbf{V}
\end{equation}

где $d_k$ --- размерность пространства ключей, а нормировка на $\sqrt{d_k}$ предотвращает чрезмерно малые градиенты при больших размерностях.

Масштабирование моделей --- увеличение числа параметров, объёма обучающих данных и вычислительных ресурсов --- привело к качественному скачку в способностях LLM. Модели семейства GPT-4, Claude 3 (Anthropic), Gemini (Google) продемонстрировали способность к выполнению сложных интеллектуальных задач, включая рассуждения, планирование, генерацию кода и решение математических задач.

\subsection{Стратегии использования LLM}

В контексте прикладных задач выделяются несколько стратегий использования LLM:

\textbf{Zero-shot prompting} предполагает подачу модели запроса без примеров выполнения задачи. Модель должна опираться исключительно на знания, полученные в процессе предварительного обучения. Этот подход прост в реализации, но часто даёт нестабильные результаты для специализированных задач.

\textbf{Few-shot prompting} \cite{brown2020gpt3} включает в запрос несколько примеров входных данных и ожидаемых ответов. Демонстрация примеров позволяет модели выявить шаблон задачи и сформировать ответ в соответствии с ожидаемым форматом. Этот подход существенно повышает качество генерации структурированных данных.

\textbf{Chain-of-Thought (CoT) prompting} \cite{wei2022chain} предполагает стимулирование модели к пошаговому рассуждению перед формированием финального ответа. В запросе явно предлагается модели «рассуждать шаг за шагом» (\textit{«think step by step»}), что приводит к существенному улучшению качества ответов на задачи, требующие многошагового логического вывода.

\textbf{ReAct (Reasoning and Acting)} \cite{yao2023react} комбинирует генерацию рассуждений с выполнением действий (вызовом инструментов), что позволяет модели взаимодействовать с внешними системами и получать дополнительную информацию в процессе решения задачи.

\textbf{Retrieval-Augmented Generation (RAG)} \cite{lewis2020retrieval} дополняет генерацию модели извлечением релевантной информации из внешних источников знаний. В контексте системной динамики RAG может использоваться для извлечения примеров моделей из базы знаний, что обогащает контекст генерации.

\subsection{Открытые языковые модели}

Наряду с проприетарными моделями активно развивается экосистема открытых LLM. Модели семейства \textbf{LLaMA} (Meta) \cite{touvron2023llama}, \textbf{Mistral} (Mistral AI) \cite{jiang2023mistral}, \textbf{Falcon} (TII), \textbf{DeepSeek} (DeepSeek AI) и другие демонстрируют качество генерации, сопоставимое с проприетарными аналогами при определённых условиях.

Открытые модели имеют ключевое преимущество для проектов, связанных с обработкой конфиденциальных данных или требующих развёртывания в закрытой инфраструктуре. В контексте проекта \texttt{sdCloud} открытые модели рассматриваются как перспективное направление для снижения зависимости от внешних API и обеспечения полного контроля над конвейером генерации.


\section{AI-ассистенты в системно-динамическом моделировании}

\subsection{Обзор существующих решений}

Интеграция больших языковых моделей с инструментами системной динамики является активно развивающимся направлением исследований. В последние годы появился ряд работ, исследующих различные аспекты этой интеграции.

Работа \cite{hu2025chatpysd} представляет систему \textbf{ChatPySD} --- плагин для ChatGPT, обеспечивающий загрузку, симуляцию и анализ моделей системной динамики непосредственно в интерфейсе чат-бота. ChatPySD использует PySD как движок симуляции и демонстрирует возможность естественно-языкового взаимодействия с существующими моделями, включая изменение параметров, запуск симуляций и визуализацию результатов. Ключевым ограничением ChatPySD является фокус на взаимодействии с готовыми моделями, а не на генерации новых моделей из текстовых описаний.

В работе \cite{akhavan2024generative} проведён обзор возможностей генеративного ИИ в контексте имитационного моделирования, включая системную динамику. Авторы выделяют три направления применения: генерация кода моделей, генерация входных данных для симуляций и генерация описаний результатов моделирования. Работа подчёркивает потенциал LLM для автоматизации рутинных задач в процессе моделирования, но отмечает необходимость человеческой верификации сгенерированных результатов.

Исследование \cite{liu2025leveraging} направлено на использование LLM для автоматической генерации CLD из текстовых описаний. Авторы демонстрируют, что GPT-4 способен извлекать переменные и причинно-следственные связи из текстовых описаний систем и формировать корректные CLD для относительно простых случаев. Качество генерации существенно зависит от сложности описываемой системы и детальности текстового описания.

\subsection{Конвейеры генерации моделей}

Совокупность существующих работ позволяет сформулировать общую архитектуру конвейера генерации системно-динамических моделей с использованием LLM. Типичный конвейер включает следующие этапы:

\begin{enumerate}
    \item \textbf{Извлечение переменных} (Variable Extraction) --- идентификация релевантных переменных из текстового описания системы.
    \item \textbf{Построение CLD} --- определение причинно-следственных связей между извлечёнными переменными и формирование петель обратной связи.
    \item \textbf{Трансляция в SFD} --- преобразование CLD в диаграмму потоков и уровней с явным выделением накопителей и потоков.
    \item \textbf{Генерация уравнений} --- формализация функциональных зависимостей в виде системы уравнений системной динамики.
    \item \textbf{Параметризация} --- назначение численных значений параметрам и начальных условий.
    \item \textbf{Верификация и валидация} --- проверка корректности сгенерированной модели и сопоставление её поведения с ожидаемым.
\end{enumerate}

Каждый этап конвейера может выполняться различными LLM с использованием специализированных промптов. Важным аспектом является итеративный характер процесса: результаты каждого этапа могут корректироваться на основе обратной связи от последующих этапов или пользователя.

\subsection{Проблемы и вызовы}

Несмотря на впечатляющие результаты отдельных исследований, область интеграции LLM и системной динамики сталкивается с рядом фундаментальных проблем:

\textbf{Галлюцинации моделей} --- LLM могут генерировать структурно корректные, но содержательно ошибочные модели. Модель может ввести переменные, не имеющие отношения к описываемой системе, или установить некорректные причинно-следственные связи.

\textbf{Сложность иерархических систем} --- для систем с большим числом переменных и сложной структурой обратных связей качество генерации существенно снижается. LLM имеют ограничения на контекст, которые могут препятствовать генерации комплексных моделей.

\textbf{Недетерминированность} --- стохастическая природа генерации LLM приводит к тому, что повторный запрос с идентичным промптом может дать различные результаты, что затрудняет воспроизводимость.

\textbf{Верификация и валидация} --- автоматическая проверка корректности сгенерированных моделей остаётся нерешённой проблемой. Требуются специализированные методы валидации графовых структур и уравнений моделей.

Проект \texttt{sdCloud} направлен на систематический подход к решению указанных проблем через внедрение многоэтапного конвейера генерации с промежуточной валидацией и использованием специализированных алгоритмов проверки топологии моделей.
